% Source: https://artofproblemsolving.com/community/c3h1530612_classic_cone_question
% Stable identity: aops-topic-1530612-q01
Let the vertex be $V$, let the base radii be $3k$ and $5k$, and let the heights be $h$ and $H$:

\begin{asy}
usepackage("amsmath");
size(4.5cm);
defaultpen(fontsize(10pt));
pair V=(0,6);
real R=2.5,r=1.5,e=0.35;
pair Y=(0,0),Z=(0,6-6*r/R);
draw(shift(Y)*scale(R,e*R)*arc((0,0),1,180,360),linewidth(0.9));
draw(shift(Y)*scale(R,e*R)*arc((0,0),1,0,180),dashed);
draw(shift(Z)*scale(r,e*r)*arc((0,0),1,180,360),blue+linewidth(0.9));
draw(shift(Z)*scale(r,e*r)*arc((0,0),1,0,180),blue+dashed);
draw((-R,0)--V--(R,0),linewidth(0.9));
draw(V--Y,dashed);
draw(Y--(R,0),dashed);draw(Z--(r,Z.y),blue+dashed);
label("$V$",V,N);
label("$3k$",(r,Z.y)+(0.12,0.12),NE,blue);label("$5k$",(R,0)+(0.12,0.12),NE);
draw((R+1.1,6)--(R+1.1,Z.y),blue,Arrows(5));label("$h$",(R+1.1,(6+Z.y)/2),E,blue);
draw((R+2.3,6)--(R+2.3,0),Arrows(5));label("$H$",(R+2.3,3),E);
draw(V--(R+2.5,6),dotted);draw((r,Z.y)--(R+1.3,Z.y),dotted+blue);draw((R,0)--(R+2.5,0),dotted);
label("\mbox{}",(min(currentpicture,user=true)+max(currentpicture,user=true))/2,black); // keeps the text after the figure black when LaTeX typesets the labels inline
\end{asy}

The cutting plane is parallel to the base, so the two right triangles formed by the axis, a radius and the slant side are similar, and the heights are in the same ratio as the radii: $h:H=3:5$. Hence the smaller cone is a scaled copy of the larger one with scale factor $\dfrac{3}{5}$, and volumes scale by the cube:
\[
\dfrac{\dfrac{1}{3}\pi(3k)^2h}{\dfrac{1}{3}\pi(5k)^2H}=\left(\dfrac{3}{5}\right)^2\cdot\dfrac{3}{5}=\boxed{\dfrac{27}{125}}.
\]
The answer $\dfrac{9}{25}$ is wrong: it uses the ratio of the base areas only, as if the two cones had the same height, but the smaller cone is also shorter by the factor $\dfrac{3}{5}$.
